\documentclass[11pt,a4paper]{article}

\usepackage[utf8]{inputenc}
\usepackage[T1]{fontenc}
\usepackage[english]{babel}
\usepackage[margin=2.5cm]{geometry}
\usepackage{booktabs}
\usepackage{graphicx}
\usepackage{amsmath}
\usepackage{xcolor}
\usepackage[colorlinks=true,linkcolor=black,citecolor=black,urlcolor=blue]{hyperref}

% Draft marker for the sections still to be written.
% Delete this macro (and every \TODO call) before handing the report in.
\newcommand{\TODO}[1]{\textcolor{red}{\textbf{[TODO:} #1\textbf{]}}}

\title{Word Embeddings and Named Entity Recognition\\
       on the QUAERO French Corpora}
\author{Alexandre Costa \and Francesc Serra}
\date{\today}

\begin{document}
\maketitle

\begin{abstract}
We build and compare word embeddings on two French corpora of very different size and domain,
and we evaluate them on medical named entity recognition (NER) together with a pretrained
Transformer. Three embedding approaches (word2vec CBOW, word2vec skip-gram and fastText CBOW)
are trained on a small medical corpus (QUAERO\_FrenchMed) and a large non-medical one
(QUAERO\_FrenchPress), compared intrinsically by cosine similarity, and then used to
initialise the embedding layer of a CNN and a BiLSTM token classifier. A fine-tuned
CamemBERT model provides the Transformer point of comparison. The Transformer reaches an
entity-level $F_1$ of $0.629$ on EMEA and $0.577$ on MEDLINE, against $0.433$ and $0.326$
for the best word-level model.
\end{abstract}

\section{Objectives}
\label{sec:objectives}

\TODO{Summarise the two lab sessions: (i) building and comparing embedding approaches on the
two corpora, (ii) evaluating them on the NER task against a Transformer.}

%% ======================================================================
\section{Session 1 --- Word embeddings}
\label{sec:tp1}

\subsection{Corpora}
\label{sec:corpora}

\TODO{Describe the two \texttt{.ospl} corpora and give their sizes (sentences, tokens, word
types). Point out that the press corpus is about $24\times$ larger than the medical one.}

\subsection{Training setup}
\label{sec:tp1-setup}

\TODO{Report the hyper-parameters required by the assignment ($\mathrm{dim}=100$,
$\mathrm{min\_count}=1$), the window size, and the number of epochs actually used. Explain why
the library default of 5 epochs had to be raised to 30.}

\subsection{Semantic similarity}
\label{sec:similarity}

\TODO{For each candidate word (\emph{patient}, \emph{traitement}, \emph{maladie},
\emph{solution}, \emph{jaune}), give the 10 nearest neighbours by cosine similarity, in the two
required comparisons: same corpus / different approaches, and same approach / different
corpora.}

\subsection{Discussion}
\label{sec:tp1-discussion}

\TODO{Interpret the tables: effect of the approach (CBOW vs skip-gram vs fastText) and effect
of the data (type and quantity).}

%% ======================================================================
\section{Session 2 --- Named entity recognition}
\label{sec:tp2}

\subsection{Dataset}
\label{sec:ner-data}

The annotated data is the QUAERO\_FrenchMed corpus in CoNLL format: five whitespace-separated
columns per line (token index, token, two unused columns, NER tag), with sentences separated by
an empty line. Tags follow the \texttt{I-TYPE} convention, with \texttt{B-TYPE} marking the
first token of an entity that immediately follows another, and \texttt{O} for tokens outside
any entity. Ten entity types occur (ANAT, CHEM, DEVI, DISO, GEOG, LIVB, OBJC, PHEN, PHYS,
PROC), giving 21 distinct tags. Two sub-corpora are provided with their own train/dev/test
splits: EMEA (drug leaflets) and MEDLINE (scientific article titles).

\TODO{Add a table of split sizes (sentences and entities per split) and note the class
imbalance: \texttt{O} covers about 79\% of tokens.}

\subsection{Adapting the provided scripts}
\label{sec:adaptation}

All three provided scripts perform \emph{sentence} classification and had to be adapted to
\emph{token} classification.

\TODO{Describe the CNN and LSTM adaptations: CoNLL loading, removing the pooling / last-time-step
reduction so that one label is predicted per token, the cross-entropy objective with
\texttt{ignore\_index} on padding, and initialising the embedding layer from the TP1 vectors.}

For \texttt{transformers\_classification.py} the changes were the following.
\texttt{AutoModelForSequenceClassification} was replaced by
\texttt{AutoModelForTokenClassification}, as required. Because BERT works on sub-word units,
the labels no longer align with the input words: the tokenizer is called with
\texttt{is\_split\_into\_words=True} and each word's tag is placed on its \emph{first}
sub-word, while continuation sub-words and the special tokens receive the label $-100$, which
the model excludes from the loss. For example \emph{Prialt} is split into
\texttt{\_Pri}/\texttt{al}/\texttt{t} and only the first piece carries \texttt{B-CHEM}. At
evaluation time the $-100$ positions are removed again, which recovers exactly one prediction
per original word, as \texttt{seqeval} requires. The objective function itself did not have to
be written: the token-classification head applies cross-entropy with
\texttt{ignore\_index}$=-100$ internally. Finally, the script's \texttt{datasets}-based CSV
preprocessing was replaced by a small \texttt{torch} \texttt{Dataset} plus
\texttt{DataCollatorForTokenClassification}, which pads each batch and pads the labels with
$-100$.

Two API names in the provided script are deprecated in \texttt{transformers}~4.57 and raise an
error: \texttt{TrainingArguments(evaluation\_strategy=...)} is now \texttt{eval\_strategy}, and
\texttt{Trainer(tokenizer=...)} is now \texttt{processing\_class}.

\subsection{Experimental setup}
\label{sec:ner-setup}

All models are evaluated with \texttt{seqeval} at the entity level, which requires both the
span and the type of an entity to be predicted exactly. We report precision, recall and $F_1$.

\TODO{Give the CNN and LSTM hyper-parameters (embedding dimension 100, filters/kernels, hidden
size, batch size, learning rate, early stopping).}

The Transformer is \texttt{camembert-base}, fine-tuned with a maximum sequence length of 128, a
batch size of 8 and a learning rate of $3\times10^{-5}$. Training runs for at most 20 epochs
with early stopping on dev $F_1$ (patience~3), and the best epoch is restored before testing.
The number of epochs matters a great deal here and is discussed in
Section~\ref{sec:q3}.

\subsection{Results}
\label{sec:ner-results}

Table~\ref{tab:results} collects the best configuration of each model on the test sets.

\begin{table}[h]
\centering
\begin{tabular}{llccc}
\toprule
Dataset & Model & Precision & Recall & $F_1$ \\
\midrule
EMEA    & CNN (press skip-gram, fine-tuned)   & 0.4623 & 0.4072 & 0.4330 \\
EMEA    & BiLSTM (press fastText, frozen)     & 0.4884 & 0.3798 & 0.4273 \\
EMEA    & \textbf{CamemBERT (fine-tuned)}     & \textbf{0.6330} & \textbf{0.6257} & \textbf{0.6293} \\
\midrule
MEDLINE & CNN (press CBOW, fine-tuned)        & ---    & ---    & 0.3244 \\
MEDLINE & BiLSTM (med fastText, fine-tuned)   & ---    & ---    & 0.3260 \\
MEDLINE & \textbf{CamemBERT (fine-tuned)}     & \textbf{0.5474} & \textbf{0.6110} & \textbf{0.5774} \\
\bottomrule
\end{tabular}
\caption{Entity-level test results (\texttt{seqeval}). Random-embedding baselines reach
$F_1=0.2648$ (CNN) and $0.2227$ (BiLSTM) on EMEA.}
\label{tab:results}
\end{table}

\TODO{Add the per-entity breakdown for the best model of each dataset, and the embedding
vocabulary coverage table.}

%% ======================================================================
\section{Answers to the questions}
\label{sec:questions}

\subsection{Which model provides the best performance on each dataset?}
\label{sec:q1}

The Transformer, on both datasets and by a wide margin: $F_1=0.629$ against $0.433$ on EMEA,
and $0.577$ against $0.326$ on MEDLINE (Table~\ref{tab:results}). The gap between the
Transformer and the word-level models is larger than any difference \emph{within} the
word-level models, whatever embeddings the latter are given.

Two secondary observations follow from the same table. The CNN and the BiLSTM are
indistinguishable from each other on both datasets, so the choice of word-level architecture
matters far less than the choice of representation. And every model scores lower on MEDLINE
than on EMEA, which is consistent with MEDLINE being made of scientific article titles: they
are shorter, denser in entities and lexically more varied than drug-leaflet sentences.

\TODO{Confirm the MEDLINE CNN/BiLSTM figures against the notebook that produced them and add
the missing precision/recall values in Table~\ref{tab:results}.}

\subsection{Differences between the embeddings learned on the two corpora}
\label{sec:q2}

\TODO{Answer from Section~\ref{sec:similarity}: the two corpora learn different senses of the
same word (the \emph{jaune} example), and the press corpus produces the better NER model
despite being out of domain because it is far larger.}

\subsection{Comparison with a Transformer model}
\label{sec:q3}

CamemBERT outperforms the best word-level model by $+0.196$ $F_1$ on EMEA ($0.6293$ against
$0.4330$) and by $+0.251$ on MEDLINE ($0.5774$ against $0.3260$), while using none of the
embeddings trained in Session~1. Three properties account for the difference.

\begin{enumerate}
  \item \textbf{Sub-word tokenisation.} There is no out-of-vocabulary case: any unseen word is
    decomposed into known pieces. The word-level models, by contrast, left between 14\% and
    42\% of their vocabulary on random vectors, depending on the source corpus.
  \item \textbf{Scale of pretraining.} CamemBERT is pretrained on billions of French tokens,
    against the $52$K and $1.25$M tokens of the two lab corpora. This is the same effect that
    made the press embeddings beat the medical ones, several orders of magnitude further along.
  \item \textbf{Contextual representations.} Each occurrence of a word receives its own vector,
    so \emph{jaune} in \emph{fièvre jaune} and in \emph{maillot jaune} are no longer forced to
    share a single representation. This removes precisely the limitation identified in
    Section~\ref{sec:q2}.
\end{enumerate}

Two caveats qualify this comparison. First, the Transformer only wins once it is trained to
convergence. A first run using the provided script's default of 5 epochs scored $F_1=0.4263$ on
EMEA --- statistically indistinguishable from the CNN --- because dev $F_1$ was still rising
steeply when training stopped: it went $0.000$, $0.000$, $0.365$, $0.471$, $0.496$ over the
five epochs and only peaked at $0.736$ on epoch~14. With 706 training sentences and a batch
size of 8, five epochs amount to just 445 optimiser steps, which is not enough for a randomly
initialised classification head over 21 labels. Comparing architectures at unequal training
budgets measures the budget rather than the architecture.

Second, the Transformer does not remove the dataset's real limitation. The rare entity types
remain unlearnable for it as well: PHEN (19 training entities) scores $0.000$, OBJC $0.133$ and
DEVI $0.171$ --- the same types on which the CNN and the BiLSTM also fail. Conversely, the
well-represented types are handled well (LIVB $0.827$, CHEM $0.725$). With 706 training
sentences spread over 10 entity types, the quantity of annotated data --- not the architecture
and not the embeddings --- is what ultimately caps every model considered here.

%% ======================================================================
\section{Conclusion}
\label{sec:conclusion}

\TODO{Wrap up: representation quality dominates architecture choice; corpus size dominates
domain match for static embeddings; a pretrained Transformer supersedes both, but the amount of
annotated data remains the binding constraint.}

%% ======================================================================
\appendix
\section{Code}
\label{sec:code}

\begin{tabular}{ll}
\toprule
File & Content \\
\midrule
\texttt{main.ipynb}            & Session 1: embedding training and semantic similarity \\
\texttt{scripts/cnn.ipynb}     & Session 2: CNN token classifier \\
\texttt{scripts/lstm.ipynb}    & Session 2: BiLSTM token classifier \\
\texttt{scripts/transformer.ipynb} & Session 2: CamemBERT token classifier \\
\bottomrule
\end{tabular}

\end{document}
